\paragraph{Result.} Under the two local estimators, each observed shortfall moves
$F$ by at most 0.52 percentage points, with joint 95\% intervals inside
$[\mn2.13, 2.13]$; the largest point bias under any estimator tried is 0.98
points (alternative clustered, low-to-center slope). For comparison, the
conditional Monte Carlo half-widths of $F$ are 2.30 to 5.24 percentage points and
the prespecified contrasts are 11.7 to 49.2 points. Estimated point contrast biases are at most
0.89 points (joint intervals inside $[\mn2.34, 3.27]$), and all six adjusted
contrasts still exclude zero at 98.333\%. The direction is not resolved: every
slope interval, for every profile and estimator, includes zero, and the point
slopes differ in sign between profiles. These estimates depend on local
linearity and the slope-transfer assumption below; they do not bound
interpolation error or establish whether matched baselines would raise or lower $F$.
The adjusted values are post hoc and descriptive and do not turn any failed check
into a pass. They adjust the 512-seed holdout $F$, so they differ from the
common-baseline $F_{\mathrm{cb}}$ by the seed-sample difference (256-seed center $F$ 66.3,
55.3 and 29.8\% against holdout $F$ 67.6, 55.9 and 27.0\% at the same $p$).

\paragraph{Assumptions.} (i) $F$ is locally linear in $\ln B$ between the
observed baseline and 3.0; only three paired points exist per profile and
curvature is unresolved. (ii) Moving $p$ is the relevant way to reach 3.0; a
baseline difference from other causes (seeds, workflow rules, admission mix)
need not carry the same change in $F$. (iii) For the alternative workflow the
original-workflow slopes are transferred; no alternative-workflow
$p$-sensitivity bank exists, and the alternative grid can check only $\gamma$.
(iv) Intervals are Monte Carlo only and exclude interpolation-model, structural
and clinical-input uncertainty. The environmental-comparison points lie on a
separate 512-seed bank with baselines 0.40 to 2.29 and serve only as secondary
slopes.

